\section{防御总纲：Defense-in-Depth 而非单点银弹}

讲者明确表示，Agent 安全没有 silver bullet，必须采用纵深防御。课程给出的防御层包括：模型加固（pre/post training）、输入护栏、输出护栏、策略执行、运行时监控、人类审批和权限管理。关键不在于每层都完美，而在于单层失守时其他层还能兜底。

\begin{knowledgebox}{典型纵深防御链路}
\begin{enumerate}
    \item \textbf{Model Hardening}：数据清洗、安全微调、对抗训练。
    \item \textbf{Input Guard}：识别并移除注入载荷、恶意上下文模式。
    \item \textbf{Planner Constraint}：限制计划生成中的高危路径。
    \item \textbf{Policy Enforcement}：执行时校验动作与权限。
    \item \textbf{Runtime Monitoring}：检测异常行为并触发回滚/暂停。
    \item \textbf{Human-in-the-Loop}：关键动作强制人工确认。
\end{enumerate}
\end{knowledgebox}

\begin{warningbox}{只做输入过滤会被自适应攻击绕过}
课程提到，攻击者会针对防御策略生成自适应样本。输入过滤只能降低低成本攻击，不足以保证高风险场景安全。必须结合策略执行与运行时约束。
\end{warningbox}

\subsection{本章小结}
真正可落地的防御不是一个模型或一个规则，而是一条多层防线。系统设计目标是把 ``单次突破'' 变成 ``连续突破''，显著抬高攻击成本。

